% Source: https://artofproblemsolving.com/community/c3h2794113_bamo8_problem_set
% Stable identity: aops-topic-2794113-q05
Let $PQ$ be a longest side and project all vertices orthogonally onto line $PQ$. Going around the polygon from $Q$ back to $P$ through the other vertices $V_1,\dots,V_k$ ($k\ge1$), suppose no $V_i$ projects onto the segment $PQ$. The projections start at $Q$ and end at $P$, so some consecutive pair of vertices on this path (possibly including $Q$ or $P$ itself) projects to points on opposite sides of segment $PQ$, at least one of them strictly outside it. The side joining them then has projection longer than $PQ$, and a segment is at least as long as its projection, contradicting the choice of $PQ$.
\begin{asy}
usepackage("amsmath");
size(4cm);
defaultpen(fontsize(10pt));
pair[] V={(0,0),(5,0),(6.5,2),(4,4),(1,3.5),(-0.8,1.8)};draw(V[0]--V[1]--V[2]--V[3]--V[4]--V[5]--cycle);draw(V[0]--V[1],linewidth(1.5));pair A=V[4];pair H=(A.x,0);draw(A--H,dashed);dot(H);label("$P$",V[0],SW);label("$Q$",V[1],SE);label("$A$",A,N);
\end{asy}

So some vertex $A=V_i$ projects onto segment $PQ$. $\blacksquare$
